\chapter*{Abstract}
\addcontentsline{toc}{chapter}{Abstract (English)}
\markboth{Abstract}{Abstract}

\label{sec:abstract-english}

A relay carries traffic between a source and a destination that have no direct path between them. This thesis asks a single question: when is a \emph{learned} relay worth using in place of a classical one, and what sets the boundary between them?

The comparison is run on one configuration, fixed in advance --- a single-relay SISO link with i.i.d.\ Rayleigh fast fading on both hops, complex baseband, Gray-coded QPSK, uncoded transmission --- leaving the relay function as the only variable. Eight strategies are evaluated, from amplify-and-forward (AF) and symbol-wise decode-and-forward (DF) to generative and sequence models, as Monte Carlo estimates with 95\% confidence intervals.

On a matched, memoryless channel the classical relay wins. Learned relays beat AF at low SNR, but from $6$~dB upward DF matches or exceeds every learned relay at no parameter cost. Capacity is not the binding constraint: a size sweep replicated over three initializations locates the smallest relay that meets the stated non-degradation budget, and \emph{channel memory} rather than parameter count sets the required context. Four parameters suffice in the tested memoryless sweep, while the evaluated three-tap channels need $73$ to $145$ and a window spanning the interference. The measured error does not rise beyond the across-initialization spread in the tested range.

The picture inverts once the channel is unknown. Against an unmodeled three-tap ISI filter both evaluated memoryless classical relays, AF and zero-threshold DF, fail outright, with DF's error rate rising as transmit power grows, while a $170$-parameter windowed network restores reliable relaying over the validated operating range and approaches the matched Viterbi reference. Two measured boundaries locate the crossover. The first is the reliability of the channel estimate: at a $10$~dB operating point, pilot-aided classical detection stays ahead at twenty pilots and loses by ten. The second is arithmetic: sequence detection costs $2M^{L}$ multiply-accumulates per symbol against the fixed relay's $2WH+4H$; the two equal near three and a half taps for the architecture deployed here and are $158\times$ apart at seven.

The evaluated learned relay is not shown to beat a model-aware matched detector. Its value in these experiments is that it still works when no such model, or no reliable estimate of one, is available. For a fixed deployed window its per-symbol arithmetic is constant; extending that window with channel memory increases the cost, while full-state sequence detection grows exponentially in the channel memory. These statements are bounded to the classical pipelines implemented here and to the impairment family represented during training.
% \begin{center}\rule{0.5\linewidth}{0.5pt}\end{center}

% \begin{flushright}
% \begin{otherlanguage}{hebrew}
% \textbf{ארכיטקטורות למידה עמוקה לתקשורת ממסר דו-קפיצתית: מחקר השוואתי של אסטרטגיות ממסר קלאסיות ומבוססות רשתות נוירונים}
% גיל צוקרמן
% חיבור זה הוגש כמילוי חלקי של הדרישות לקבלת תואר מגיסטר למדעים (M.Sc.)
% בית הספר להנדסת חשמל — אוניברסיטת תל אביב
% 2026
% \end{otherlanguage}
% \end{flushright}

\clearpage

%% ── Acknowledgments ──────────────────────────────────────────
\cleardoublepage
\phantomsection
\chapter*{Acknowledgments}
\addcontentsline{toc}{chapter}{Acknowledgments}
\markboth{Acknowledgments}{Acknowledgments}

I would like to thank my supervisors, Dr. Anatoly Khina and Prof. Raja Giryes at Tel Aviv University, for their guidance and support throughout this research.
This work was conducted in the School of Electrical and Electronic Engineering, Tel Aviv University.
